% !TEX root = ../main.tex
\chapter{Course Logistics}\label{app:logistics}
\begin{paracol}{2}

This appendix collects the practical, non-mathematical information about how the two courses were
run: syllabus and grading, calendar structure, problem-set and assignment logistics, the 2024
investment game, the 2024 group project and final paper, and the 2013 instructor interview. It is
assembled from the \file{pages/} subdirectories of the two OCW course sites (\file{syllabus},
\file{calendar}, \file{assignments}, \file{problem-sets}, \file{investment-game},
\file{group-project}, \file{final-paper} for 2024; \file{syllabus}, \file{instructor-insights},
\file{projects}, \file{related-resources} for 2013), reporting only what those pages state.

\section{18.642, Fall 2024}

\subsection{Syllabus and grading}

Two 1.5-hour lectures per week. Prerequisites: 18.03 (Differential Equations), 18.05 or 18.600
(Probability), and 18.06 (Linear Algebra); prior finance knowledge is not required. Mathematics
lectures were given by MIT mathematicians (mainly Peter Kempthorne), applications lectures by
industry professionals (a mix of investment banks, FinTech and hedge funds).

The class had five problem sets, an optional investment game, a mid-semester group project (a
written lecture note plus an in-class presentation), and, in place of a final exam, a final paper on
a math finance topic of the student's choice. The grading page states the weights explicitly:
\begin{center}
\begin{tabular}{@{}ll@{}}
\toprule
Component & Weight\\
\midrule
Homework (problem sets) & 40\%\\
Group project presentation & 20\%\\
Participation & 10\%\\
Final paper & 30\%\\
\bottomrule
\end{tabular}
\end{center}

\subsection{Calendar structure}

The calendar page lists 26 lecture slots, alternating math lectures (mostly Kempthorne, covering the
topics of \cref{part:2024}) with applications lectures by outside speakers (e.g.\ Jeff Shen of
BlackRock on quantitative equity investing, L3; Andrew Gunstensen of Mizuho on rates, L7; Stefan
Andreev of Two Sigma on PCA, L9; Andrew W.~Lo on AI in biomedical portfolios, L18; John Hull on
machine learning, L23). Problem sets were due at lectures 3, 7, 9, 18 and 22. Lectures 15--17 were
student presentations of the group project, and lecture 26 was the final paper presentations.

\subsection{Problem sets}

Five problem sets, each covering the math lecture topics up to that point (see
\cref{app:materials} for the file-to-chapter mapping). The problem-sets page lists two accompanying
readings: Gron, ``Optimal Portfolio Choice and Stochastic Volatility'' (\emph{Applied Stochastic
Models in Business and Industry}, 2012), referenced by Problem 2 of Problem Set~2, and, for the PCA
problem, a note on principal component analysis of interest rates plus sample R code (both used in
\cref{ch:pca}); Problem Set~5 references Parkinson, ``The Extreme Value Method for Estimating the
Variance of the Rate of Return'' (\emph{The Journal of Business}, 1980) and ships a companion R
Markdown volatility case study.

\subsection{Investment game}\label{appC:game}

The investment game (ungraded) is stated in full detail in \cref{ex:game} of \cref{ch:markets}; it
is not repeated here. Two points from the course page are purely administrative and not part of the
exercise: students had to submit their choice of instrument to the instructors before 14 September
2024, and their optional switch before 12 October 2024, and at the end of the term the instructors
announced the top three performers by sharing the tracking spreadsheets.

\subsection{Group project}

Groups worked on an introductory, written lecture note aimed at fellow students with no prior
knowledge of the topic; each group member submitted their own 5--7 page draft on their section. The
first submission (``Draft 1'') was graded only on evidence of individual, focused effort (5+ pages
of draft text), and closed with a reflection paragraph asking for specific feedback. It was followed
by a peer review of another group's draft, revision, and a 15--20 minute oral presentation in class
(lectures 15--17 on the calendar).

The topics offered in 2024 were: How Human Fallacies Lead to Market Anomalies; Regression Modeling
of UK Property Prices; The Use of Machine Learning in Predicting Future Interest Rates; Alternative
Data; Jump-Diffusion Model; Maximizing Utility While Minimizing Risk; Fourier Analysis and Its
Applications in the Stock Market; Black Scholes Merton Model; Implied Volatility; Partial
Differential Equation Models for Option Pricing; Manipulating the Sports Betting Market; Time
Series Modeling; Analyzing Volatility Curves in Option Pricing; The Black-Scholes Model; and
Delta-Hedging.

\subsection{Final paper}

Every student wrote their own final paper, individually graded even when it continued a group
project topic; sections of a joint paper had to be clearly attributed to their author. Papers were
due on the last day of lectures, typically 10--15 pages, graded in part on mathematical content,
writing quality, and organization. As with the group project, a Draft~1 was required first (3--4
polished pages plus 1--2 pages of motivation, or 5--7 pages of developing notes, with a closing
reflection section), assessed only on evidence of focused effort, followed by peer review the next
week.

The final-paper page lists 33 topics proposed for 2024, spanning market microstructure (credit
expansion and land prices, election effects on volatility), derivatives pricing (SABR, Heston,
jump-diffusion Monte Carlo, comparative Black-Scholes/Monte Carlo studies), portfolio and risk
methods (Black-Litterman, PCA, tail-risk hedging), machine learning applications (LSTM volatility
forecasting, reinforcement learning for Greeks hedging, AI-driven trading), and a few lighter
entries (sports-betting binary options, ``The Wolf of Wall Street'' and stock-market cinematography,
DEI-focused ETFs). The full list is not reproduced here; consult the source page for all 33 titles.

\section{18.S096, Fall 2013}

\subsection{Syllabus, goals and grading}

Same meeting pattern (two 1.5-hour lectures/week) and a longer prerequisite list: 18.01, 18.02,
18.03, 18.05 or 18.440, and 18.06. The course had no exams: problem sets were due two weeks after
each math lecture, and the final grade was 75\% homework, 25\% final paper. The syllabus page lists
ten explicit course goals, later repeated verbatim on the instructor-insights page, e.g.\ deriving
the price-yield relationship and convexity, bootstrapping a yield curve, computing standard VaR,
estimating option volatility, deriving Black--Scholes via risk-neutral arguments, PCA as matrix
decomposition, limiting theorems, It\^o's lemma, and Girsanov's theorem and change of measure --
i.e.\ essentially the syllabus of \cref{part:2024,part:2013} combined. The course also offered an
optional field trip to Morgan Stanley's New York offices.

\subsection{Instructor insights (2013)}

The course grew out of Jake Xia's and Vasily Strela's wish, as MIT alumni working in industry, to
give back to the department; it was first taught in Fall 2012 as invited industry lectures only, and
the pure-math component was added the following year to better integrate theory and practice.
Enrollment was about 50 students, mostly undergraduates with a few graduate students, mainly
mathematics majors with some from EECS, economics, and aero/astro; about half had no prior finance
background. Students were expected to spend about 12 hours/week on the course, split between the two
lecture types and, outside class, problem sets (due two weeks after each lecture) and the final
paper.

\begin{coursenote}
The instructors' own reflection: designing problem sets for a mixed-background cohort was an
acknowledged difficulty (``Future offerings will distinguish core versus advanced problems''), which
is one motivation, together with the addition of the applications lectures, for how the course
evolved from 2013 to the 2024 edition mapped in \cref{app:materials}.
\end{coursenote}

\subsection{Final paper topics (2013)}\label{appC:2013paper}

Twenty-five percent of the grade was the final paper (see above); the projects page gives sample
topics -- portfolio risk management, regime-shift modeling, low-volatility investing, modeling
financial bubbles (citing Didier Sornette's work), the Black--Scholes/heat-equation correspondence,
hybrid FX/rate products, HJM versus short-rate models, and Ross recovery -- plus a list of topics
actually chosen by students the previous year, mostly variations on the same themes (Black-Scholes
to heat-equation derivations, HJM/Ho-Lee models, FX jump-diffusion bond pricing, Asian options,
minimal-entropy martingale measures, PCA of energy/currency swap curves around the 2008 crisis, and
Monte Carlo/Heston option pricing).

\subsection{Related resources (2013)}

The related-resources page lists six textbooks (Hull's \emph{Options, Futures, and Other
Derivatives}; Baxter and Rennie's \emph{Financial Calculus}; Wilmott, Howison and Dewynne's
\emph{The Mathematics of Financial Derivatives}; Grinold and Kahn's \emph{Active Portfolio
Management}; Fabozzi, Focardi and Kolm's \emph{Financial Modeling of the Equity Market}; and Tsay's
\emph{Analysis of Financial Time Series}) together with a financial glossary of about 40 terms
(alpha, beta, the Greeks, VaR, swaptions, zero-coupon bonds, and so on). The glossary substantially
overlaps with the terminology already introduced across \cref{part:2024,part:2013}; \cref{ex:glossary}
in \cref{ch:markets} is the course's own suggested exercise for building such a glossary, so it is
not reproduced again here.

\section{Sources not covered}

The 2024 \file{pages/week-1} through \file{week-15} pages are a lecture-by-lecture index used to
build \cref{app:materials} and are not repeated here. No \file{case-studies} or
\file{lecture-notes} index page beyond what is already listed in \cref{app:materials} added further
logistics content for the 2013 course.
